\section{An MPI determinant}
\label{sec:mpidesign}

The second part of the assignment computes the determinants with MPI~\cite{mpi40} on
several machines and compares the result with the cached sequential program.

\subsection{The elimination as a PLU factorization}

Gaussian elimination with row exchanges factors a matrix as
$PA = LU$~\cite{golub2013matrix}. At step $k$ it picks a nonzero pivot in column $k$, swaps
its row into place, divides the entries below the pivot by it to get the multipliers
$\ell_{ik}$, and subtracts $\ell_{ik}$ times row $k$ from every row $i$ below. The
determinant is the product of the pivots, with the sign flipped once for every row
exchange. If a column has no nonzero entry at or below the diagonal, the determinant is
zero. Over a finite field there is no rounding, so any nonzero entry is as good a pivot as
any other, and we take the first one.

\subsection{Distributing by columns}

The assignment text asks us to watch how pivot selection and matrix updates affect network
communication. Step $k$ needs two things: the pivot, and the multipliers below it. Both are
in column $k$. We therefore give every process whole columns, so that the process that owns
column $k$ can find the pivot and compute the multipliers from its own memory, without
asking anyone. It then tells the others the pivot row and the multipliers, and each process
updates the columns it owns: for a column $j$ it swaps the two rows, reads $x = a_{kj}$, and
computes $a_{ij} \leftarrow a_{ij} - \ell_{ik}\,x$ for the rows below.

The alternative is to distribute rows. Then the entries of column $k$ are spread over all
processes. Finding a nonzero pivot in general takes a reduction over all of them at every
step, and a row exchange can mean two processes trading rows. With columns, a row exchange
is something every process does locally to its own columns once it knows the pivot row.

\subsection{Panels}

Sending one message per column would cost $N$ broadcasts. We group $b$ neighbouring columns
into a panel and deal the panels out to the $p$ processes in turn, a one-dimensional form
of the block-cyclic layout used by ScaLAPACK~\cite{blackford1997scalapack}. Dealing them in turn keeps the load even,
because every process still owns columns near the end of the matrix. The owner of a panel
works through its $b$ columns one at a time. For each it applies the earlier steps of the
same panel, finds the pivot and computes the multipliers. It then broadcasts the whole
panel at once: a flag for a zero determinant, the $b$ pivot rows and the $b$ columns of
multipliers. Listing~\ref{lst:apply} is the routine every process runs on each of its
remaining columns.

\begin{lstlisting}[caption={Applying one panel to one column. \code{first} is the panel's
first column.},label={lst:apply}]
for (step = 0; step < steps; ++step) {
  k = first + step;
  if (pivots[step] != k)
    swap(column[k], column[pivots[step]]);
  x = column[k];
  if (x.is_zero()) continue;
  m = multipliers + step * rows;
  for (i = k + 1; i < n; ++i)
    column[i] = column[i].add(
      x.multiply(m[i - first]).negate());
}
\end{lstlisting}

The inner loop performs the same three field operations per entry as the sequential
program: a multiplication, a negation and an addition. The steps of a panel are applied to
a column one after another, in order, which gives the same result as the row-by-row
elimination.

\subsection{What travels over the network}

Rank~0 generates the matrix and sends each process its columns once. After that the only
messages are the panel broadcasts, $\lceil N/b \rceil$ of them. Together they carry each
multiplier once, about $N^2/2$ field elements whatever the value of $b$, so $b$ trades the
number of messages against their size and against the time the other processes wait while
one process prepares a panel.

The cached field matters here as well. A cached element is a plain 16-bit value, so a
column or a panel is sent as raw bytes, two per element. An element of the original library
owns heap memory and would have to be packed and unpacked. Every process keeps its own
tables. Before timing, each process fills them, which took \CacheWarmSeconds~s per process
on \MpiMachine.

\subsection{Checking the MPI program}

In every run rank~0 also computes the determinant with the sequential elimination on the
same matrix, and the program exits with an error if the two results differ. A test script
runs \MpiTestCount{} configurations, \MpiTestVerdict. They cover one to four processes,
panel widths that do and do not divide $N$, matrices smaller than the number of processes,
sparse matrices, permutation matrices and matrices with a zero leading column, which force
row exchanges, and singular matrices. For $N = 100$, 200 and 400 the script also compares
the results with the determinants printed by the original assignment program.
